\begin{table}[H]
\centering
\small
\caption{Dominância geográfica e diversidade nas 12 redes individuais}
\label{tab:geographic_summary}
\begin{tabular}{llrrrcc}
\hline
\textbf{Rede} & \textbf{Países} & \textbf{Medalhas} & \textbf{Top 3 (\%)} & \textbf{Entropia} & \textbf{Gini} & \textbf{\#1 País} \\
\hline
Futebol M & 34 & 1335 & 21.6 & 4.81 & 0.350 & BRA \\
Futebol F & 8 & 368 & 61.7 & 2.77 & 0.317 & USA \\
Basquete M & 18 & 723 & 55.9 & 3.39 & 0.528 & USA \\
Basquete F & 14 & 429 & 52.9 & 3.32 & 0.408 & USA \\
Atletismo 100m M & 18 & 88 & 64.8 & 2.94 & 0.626 & USA \\
Atletismo 100m F & 17 & 63 & 58.7 & 3.29 & 0.543 & USA \\
Natação 100m Livre M & 20 & 81 & 59.3 & 3.24 & 0.601 & USA \\
Natação 100m Livre F & 15 & 73 & 61.6 & 3.05 & 0.559 & USA \\
Judô -90kg M & 24 & 52 & 36.5 & 4.20 & 0.378 & JPN \\
Judô -70kg F & 12 & 28 & 42.9 & 3.38 & 0.292 & FRA \\
Boxe Welter M & 35 & 89 & 22.5 & 4.82 & 0.356 & USA \\
Boxe Médio F & 5 & 8 & 75.0 & 2.25 & 0.150 & CHN \\
\hline
\multicolumn{7}{l}{\footnotesize Top 3: \% de medalhas concentradas nos 3 países dominantes} \\
\multicolumn{7}{l}{\footnotesize Entropia: diversidade geográfica (maior = mais diverso)} \\
\multicolumn{7}{l}{\footnotesize Gini: desigualdade de distribuição (0 = igualitário, 1 = concentrado)} \\
\end{tabular}
\end{table}
